\section{Which proof ingredients can actually transfer?}
\label{sec:architecture}

\subsection{One arithmetic object, two incompatible-looking bounds}

The core construction produces algebraic modular forms and chains of
theta lifts with integral structures. Divisibility at finite places
reduces their effective size, while archimedean estimates control their
norms. A second-moment argument ensures that enough of the resulting
vectors are independent. A separate argument using determinants,
congruences, local Galois cohomology, and auxiliary twists gives a
lower bound on the \emph{same} corrected covolume.

The mathematical point is the cancellation of matching main terms.
Writing \(D\) for target dimension, \(n\) for the number of independent
vectors, \(V\) for a logarithmic squared covolume bound, and \(\delta\)
for the index loss, the comparison has the schematic form
\[
 H-\mathrm{err}_{\rm low}
 \ \leq\ \frac{V+2\delta}{D}
 \ \leq\ H-\frac nD\,mR+\mathrm{err}_{\rm up}.
\]
Subtracting \(H\) leaves an upper bound for a weighted prime remainder
\(R\). In the actual source, the result is
\begin{equation}
\label{eq:star}
 m\,\mathrm{totalRemainder}\leq f a(\log b)^5.
\end{equation}
This is \texttt{star\_inequality}.\cite{mainineq}

\begin{center}
\small
\begin{tabular}{p{0.19\linewidth}p{0.34\linewidth}p{0.37\linewidth}}
\toprule
Input & Mathematical role & What a new target must supply\\
\midrule
Short forms & Divisibility and norm estimates for a spanning source
family & A local divisibility mechanism encoding its prime weights\\[3pt]
Enough forms & Near-full target span at controlled translate cost
& A quantitative second moment with compatible integral structures\\[3pt]
Long integral forms & Lower bound for covolume plus index
& Determinant identities, manageable congruence denominators,
and a common height bound\\[3pt]
Analytic step & One-sided remainder bounds exclude zeros
& The correct Euler product and control of possible pole cancellation\\
\bottomrule
\end{tabular}
\end{center}

The source names the first three stages Theorems A, B, and C.
Theorem B supplies a proportion at least \(1-4/b^2\) of the target
dimension. A fixed positive proportion is not interchangeable with
that statement, because the vectors used to complete a basis can be
expensive.\cite{shortenough,longforms}

\subsection{Quantifiers that must survive a transfer}

The lower bound does not require one globally favorable twist at every
coefficient prime. It permits a different twist at each \(p\)-adic
embedding. However, the determinant ratios for the selected family must
have a \emph{compatible height bound}. Likewise, the integral matrices
entering the index estimate must work uniformly over the allowed twists
at that place. Local existence without these simultaneous compatibility
conditions is insufficient.\cite{twistedlattice,heights}

The source's abstract product-formula lemma gives, schematically,
\[
 0\leq c+2\delta+
 n'\bigl(b_{\rm auxiliary}+b_{\rm scalar}\bigr)+b_{\rm determinant}.
\]
Here \(n'\) counts completion vectors, \(c\) bounds the covolume, and the
other terms pay for their lengths and for heights.
The exact interfaces are already useful outside the final Dirichlet
application. What is difficult is constructing a new family for which
all four terms match with the required asymptotic precision.

The Selmer quantity refined in this report belongs to the denominator
part of this comparison. Improving it is therefore relevant even though
the global prime-sum comparison is not yet available for a new target.
No statement below asserts that removing a rank factor in this one
term removes the rank factors in the entire proof.

\subsection{Why the final analytic criterion is the easier part}

The final scale selection takes the number of lifts \(m\) on the order
of \(\sqrt{x}\) while the logarithmic losses grow slower than arbitrarily
small positive powers of \(x\). The result is an eventual one-sided
remainder estimate at square-root scale. The zero-exclusion step then
uses a transform with prescribed poles and a favorable quadratic
twist for the zero under consideration.\cite{onesided,criterion}

This explains why generalizations to other prime-weight functions are
plausible. It also explains why writing such a criterion alone is not a
solution. A nonabelian character can be inserted into a formal Euler
sum, but its arithmetic divisibility interpretation and determinant
comparison must still be proved. The supplied corpus already contains
many such analytic reductions.

\subsection{Other independently reusable upstream arithmetic}

The background and internal lemmas include more than zero detectors.
For example, a local Kummer-chart theorem bounds a crystalline
orthogonality defect by
\[
 11f\bigl[m(a+1)+\operatorname{rank}(L)(W+1)\bigr]
       \operatorname{length}_O(O/pO).
\]
The parameter \(m\) counts low Hodge--Tate jumps; it is not automatically
the rank. Its hypotheses include an unramified Kummer chart, opposite
weight windows, a large gap, and a restriction excluding nearby
multiples. A companion result makes the defect zero under stronger
window assumptions.\cite{localdefect}
This is a local input potentially relevant to other integral Selmer
comparisons. It is not a proof of global Bloch--Kato vanishing.
The attached source notes retain its exact hypotheses instead of
promoting its conclusion to a more general statement.
